\subsection{Thiết kế Kafka cho Event-Driven Architecture và Event Sourcing}

Trong hệ thống thương mại điện tử này, Kafka được sử dụng như một xương sống giao tiếp bất đồng bộ giữa các microservice. Mô hình này giúp các service hoạt động độc lập, giảm phụ thuộc trực tiếp qua HTTP, đồng thời hỗ trợ khả năng mở rộng, tính chống chịu lỗi, và khả năng ghi nhận lịch sử sự kiện (event sourcing).

\subsubsection{Vai trò của Kafka trong kiến trúc}

Kafka đảm nhận ba chức năng chính:

\textbf{1.1. Event Streaming Layer – Lớp luồng sự kiện}

Các microservice không gọi trực tiếp sang nhau khi xử lý nghiệp vụ phức tạp trừ khi chúng nằm trong luồng giao dịch phân tán (Saga). Thay vào đó:

\begin{itemize}
    \item Service phát ra event (Producer)
    \item Các service khác lắng nghe topic liên quan (Consumer)
    \item Mỗi service xử lý phần việc của mình mà không ảnh hưởng thời gian phản hồi API chính
\end{itemize}

Ngoại lệ là giao tiếp đồng bộ (Sync) qua Internal API được sử dụng khi cần orchestrator service điều phối một giao dịch phân tán (Saga) để đảm bảo tính nhất quán (ví dụ: quy trình đặt hàng).

Điều này giúp frontend giao tiếp nhẹ nhàng, còn backend xử lý nội bộ qua Kafka.

\textbf{1.2. Event Sourcing – Lưu vết các sự kiện quan trọng}
Một số sự kiện trong hệ thống được xem như "nguồn sự thật" (source of truth) cho các chức năng khác, ví dụ:

Các sự kiện này có vai trò nền tảng khi được tiêu thụ bởi nhiều service khác nhau để thực thi các logic nghiệp vụ tiếp theo, hoặc được dùng để tái dựng lại trạng thái dữ liệu và tạo các báo cáo tổng hợp.

Tuy nhiên, Event Sourcing không áp dụng cho các API đặc biệt như Google Login ngay từ ban đầu, do quá trình này không chủ động tạo tài khoản, thông tin người dùng được cung cấp từ bên thứ ba thường không đầy đủ, và quá trình đăng nhập diễn ra bên ngoài hệ thống. Mặc dù vậy, hệ thống vẫn sẽ phát sự kiện use\_registered sau lần đăng nhập Google đầu tiên khi tài khoản được tạo nội bộ.

\textbf{1.3. Integration Layer – Kết nối giữa các service qua Kafka}

Trong kiến trúc của hệ thống, các service được tích hợp với nhau thông qua Kafka, đóng vai trò là kênh sự kiện trung gian.
Thay vì gọi trực tiếp (synchronous API call), các service giao tiếp bằng cách phát (publish) và lắng nghe (subscribe) các sự kiện nghiệp vụ. Điều này giúp giảm phụ thuộc giữa các service và hỗ trợ mở rộng theo hướng event-driven.

\subsubsection{Các topic Kafka tiêu biểu trong hệ thống}

Để hỗ trợ kiến trúc microservice theo hướng sự kiện (event-driven), hệ thống sử dụng một số nhóm topic Kafka tiêu biểu. Các topic này giúp truyền tải thay đổi trạng thái và dữ liệu nghiệp vụ giữa các service mà không cần phụ thuộc trực tiếp vào nhau.

\textbf{a/ User Topics}
\begin{itemize}
    \item \texttt{user.register.local}: Phát sinh khi người dùng tạo tài khoản mới.
    \item \texttt{user.login.local}: Ghi nhận sự kiện đăng nhập thành công.
    \item \texttt{user.updated}: Cập nhật thông tin người dùng.
    \item \texttt{user.behavior.tracked}: Lưu vết hành vi để phục vụ gợi ý sản phẩm/AI.
\end{itemize}

\textbf{b/ Cart Topics}
\begin{itemize}
    \item \texttt{cart.checked\_out}: Người dùng xác nhận thanh toán; Order Service sẽ lắng nghe để tạo đơn hàng.
    \item \texttt{cart.updated}: Ghi nhận thay đổi giỏ hàng, phục vụ phân tích hành vi.
\end{itemize}

\textbf{c/ Order Topics}
\begin{itemize}
    \item \texttt{order.created}: Đơn hàng mới được tạo từ quá trình checkout.
    \item \texttt{order.status.updated}: Cập nhật trạng thái đơn hàng.
    \item \texttt{order.completed}: Đơn hàng hoàn tất; Product và Admin Service sử dụng để thống kê.
    \item \texttt{order.saga.failed}: Ghi nhận khi một Saga Đặt hàng bị hủy do lỗi (compensation).
\end{itemize}

\textbf{d/ Product Topics}
\begin{itemize}
    \item \texttt{product.created}: Sản phẩm mới được thêm vào hệ thống.
    \item \texttt{product.updated}: Sản phẩm được chỉnh sửa bởi admin.
    \item \texttt{product.sales.updated}: Ghi nhận dữ liệu bán hàng sau mỗi đơn thành công.
\end{itemize}

\textbf{e/ Admin Topics}
\begin{itemize}
    \item \texttt{admin.log.created}: Lưu lại hành động của quản trị viên.
    \item \texttt{admin.report.requested}: Các yêu cầu tạo báo cáo thống kê.
    \item \texttt{admin.report.generated}: Báo cáo hoàn thành và ghi vào database.
\end{itemize}

\textbf{f/ AI Topics}
\begin{itemize}
    \item \texttt{ai.recommendation.requested}: Yêu cầu tạo gợi ý sản phẩm dựa trên hồ sơ hoặc hành vi.
    \item \texttt{ai.recommendation.generated}: Kết quả gợi ý sẵn sàng để trả về frontend.
    \item \texttt{ai.user\_behavior.received}: AI nhận dữ liệu hành vi phục vụ huấn luyện mô hình.
\end{itemize}

\subsubsection{Thiết kế sự kiện (Event-Sourcing)}

Trong mô hình Event-Sourcing, mỗi service ghi lại các thay đổi trạng thái dưới dạng sự kiện (event) và xuất bản chúng lên Kafka. 
Kafka đóng vai trò như một \textit{event store phân tán}, cho phép các service khác nhau tiêu thụ sự kiện để cập nhật trạng thái cục bộ hoặc kích hoạt luồng nghiệp vụ liên quan. 
Dưới đây là danh sách các topic và sự kiện tiêu biểu theo từng service.

\vspace{0.5em}

\textbf{a/ User Service} \\
\textbf{Topic:} \texttt{user-events}
\begin{itemize}
    \item UserRegistered
    \item UserUpdatedProfile
    \item UserLoggedIn
    \item UserLoggedOut
    \item UserDeleted
\end{itemize}

\textbf{b/ Product Service} \\
\textbf{Topic:} \texttt{product-events}
\begin{itemize}
    \item ProductCreated
    \item ProductUpdated
    \item ProductPriceChanged
    \item ProductInventoryAdjusted
    \item ProductDisabled
\end{itemize}

\textbf{c/ Cart Service} \\
\textbf{Topic:} \texttt{cart-events}
\begin{itemize}
    \item CartCreated
    \item ItemAddedToCart
    \item ItemRemovedFromCart
    \item ItemQuantityChanged
    \item CartCheckedOut
\end{itemize}

\textbf{d/ Order Service} \\
\textbf{Topic:} \texttt{order-events}
\begin{itemize}
    \item OrderCreated
    \item OrderPaid
    \item OrderCancelled
    \item OrderDelivered
    \item OrderFailed
    \item OrderCompensationReverted
\end{itemize}

\textbf{f/ Admin Service} \\
\textbf{Topic:} \texttt{admin-events}
\begin{itemize}
    \item AdminLoggedAction
    \item ReportRequested
    \item ReportGenerated
\end{itemize}

\textbf{g/ AI Service} \\
\textbf{Topic:} \texttt{ai-events}
\begin{itemize}
    \item UserBehaviorTracked
    \item ProductViewed
    \item ModelTrainingTriggered
    \item RecommendationGenerated
\end{itemize}















